% Định dạng theo mẫu đồ án tốt nghiệp của Trường CNTT&TT, ĐHBK Hà Nội (bản LaTeX).
% Biên dịch bằng XeLaTeX (Tectonic): Times New Roman 13pt, giãn dòng 1,2, lề trái 3,5 cm.
\usepackage[fontsize=13pt]{scrextend}
\usepackage{fontspec}
\setmainfont{Times New Roman}
\setmonofont[Scale=0.84]{Menlo}
\usepackage{polyglossia}
\setdefaultlanguage{vietnamese}
% polyglossia đặt lại tên mục khi vào \begin{document}, nên ghi đè trong \captionsvietnamese
\gappto\captionsvietnamese{%
  \renewcommand{\contentsname}{MỤC LỤC}\renewcommand{\listfigurename}{DANH MỤC HÌNH VẼ}%
  \renewcommand{\listtablename}{DANH MỤC BẢNG BIỂU}\renewcommand{\bibname}{TÀI LIỆU THAM KHẢO}%
  \renewcommand{\figurename}{Hình}\renewcommand{\tablename}{Bảng}}

\usepackage[a4paper,left=3.5cm,right=2.5cm,top=3cm,bottom=3cm,headheight=16pt]{geometry}
\usepackage{setspace}
\setstretch{1.2}
\setlength{\parindent}{0.75cm}
\setlength{\parskip}{4pt}

\usepackage{amsmath}
\usepackage{unicode-math}
\setmathfont{texgyretermes-math.otf}
\usepackage{graphicx}
\usepackage{float}
\usepackage{pdflscape}
\usepackage{array,multirow,longtable,tabularx,makecell}
\usepackage{xcolor}
\usepackage{enumitem}
\setlist{itemsep=2pt,topsep=4pt,parsep=0pt}
\setlist[itemize]{label=--,leftmargin=1.3cm}
\setlist[enumerate]{leftmargin=1.3cm}
\usepackage{listings}
\usepackage{algorithm}
\usepackage[noend]{algpseudocode}
\floatname{algorithm}{Thuật toán}
\renewcommand{\algorithmicrequire}{\textbf{Đầu vào:}}
\renewcommand{\algorithmicensure}{\textbf{Đầu ra:}}
\renewcommand{\algorithmicif}{\textbf{nếu}}
\renewcommand{\algorithmicthen}{\textbf{thì}}
\renewcommand{\algorithmicelse}{\textbf{ngược lại}}
\renewcommand{\algorithmicfor}{\textbf{với mỗi}}
\renewcommand{\algorithmicdo}{\textbf{làm}}
\renewcommand{\algorithmicreturn}{\textbf{trả về}}
\renewcommand{\algorithmicprocedure}{\textbf{thủ tục}}
\usepackage{pgfplots}
\pgfplotsset{compat=1.18}
\usepackage{caption}
\captionsetup{labelsep=colon,font=small,labelfont=bf,skip=6pt}
\captionsetup[table]{position=below}
\usepackage{titlesec}
\usepackage{titletoc}
\usepackage{fancyhdr}
\usepackage{tocloft}
\usepackage[hidelinks]{hyperref}
\usepackage{indentfirst}

% Màu dùng chung
\definecolor{codebg}{HTML}{F4F5F6}
\definecolor{coderule}{HTML}{D7DADF}
\definecolor{codekw}{HTML}{1D4F91}
\definecolor{codestr}{HTML}{7A3E00}
\definecolor{codecmt}{HTML}{5D636A}

% Tên các mục
\renewcommand{\contentsname}{MỤC LỤC}
\renewcommand{\listfigurename}{DANH MỤC HÌNH VẼ}
\renewcommand{\listtablename}{DANH MỤC BẢNG BIỂU}
\renewcommand{\bibname}{TÀI LIỆU THAM KHẢO}
\renewcommand{\figurename}{Hình}
\renewcommand{\tablename}{Bảng}
\renewcommand{\lstlistingname}{Mã nguồn}

% Tiêu đề chương: "CHƯƠNG 1. GIỚI THIỆU ĐỀ TÀI", in hoa, căn giữa
\titleformat{\chapter}[block]{\centering\bfseries\fontsize{14}{18}\selectfont}{CHƯƠNG \thechapter.}{0.35em}{\MakeUppercase}
\titlespacing*{\chapter}{0pt}{-12pt}{22pt}
\titleformat{\section}{\bfseries\fontsize{13}{16}\selectfont}{\thesection}{0.9em}{}
\titlespacing*{\section}{0pt}{10pt}{4pt}
\titleformat{\subsection}{\bfseries\fontsize{13}{16}\selectfont}{\thesubsection}{0.9em}{}
\titlespacing*{\subsection}{0.6cm}{8pt}{3pt}
\titleformat{\subsubsection}{\bfseries\itshape\fontsize{13}{16}\selectfont}{\thesubsubsection}{0.9em}{}
\titlespacing*{\subsubsection}{0.6cm}{6pt}{2pt}
\setcounter{secnumdepth}{2}
\setcounter{tocdepth}{2}

% Mục lục: "CHƯƠNG 1. ..." ở cấp chương
\renewcommand{\cftchappresnum}{CHƯƠNG }
\renewcommand{\cftchapaftersnum}{.}
\setlength{\cftchapnumwidth}{6.6em}
\renewcommand{\cftchapfont}{\bfseries}
\renewcommand{\cftchappagefont}{\bfseries}
\renewcommand{\cftchapleader}{\cftdotfill{\cftdotsep}}
\setlength{\cftbeforechapskip}{6pt}
\setlength{\cftsecindent}{0pt}
\setlength{\cftsecnumwidth}{2.6em}
\setlength{\cftsubsecindent}{2.6em}
\setlength{\cftsubsecnumwidth}{3.4em}
\renewcommand{\cfttoctitlefont}{\hfill\bfseries\fontsize{14}{18}\selectfont}
\renewcommand{\cftaftertoctitle}{\hfill}
\renewcommand{\cftloftitlefont}{\hfill\bfseries\fontsize{14}{18}\selectfont}
\renewcommand{\cftafterloftitle}{\hfill}
\renewcommand{\cftlottitlefont}{\hfill\bfseries\fontsize{14}{18}\selectfont}
\renewcommand{\cftafterlottitle}{\hfill}
\renewcommand{\cftfigpresnum}{Hình~}
\renewcommand{\cfttabpresnum}{Bảng~}
\setlength{\cftfignumwidth}{5.2em}
\setlength{\cfttabnumwidth}{5.2em}
\setlength{\cftfigindent}{0pt}
\setlength{\cfttabindent}{0pt}

% Đầu trang: tên chương và đường kẻ; chân trang: số trang ở giữa
\pagestyle{fancy}
\fancyhf{}
\renewcommand{\chaptermark}[1]{\markboth{CHƯƠNG \thechapter.\ \MakeUppercase{#1}}{}}
\fancyhead[L]{\small\leftmark}
\fancyfoot[C]{\thepage}
\renewcommand{\headrulewidth}{0.4pt}
% Trang mở chương dùng kiểu plain mặc định (chỉ số trang ở giữa); định nghĩa lại plain bằng fancyhdr sẽ xóa đầu trang của các trang sau.

% Mã nguồn
\lstset{
  basicstyle=\ttfamily\small\setstretch{1.0},
  backgroundcolor=\color{codebg},
  frame=single, rulecolor=\color{coderule}, framesep=4pt,
  keywordstyle=\color{codekw}\bfseries, stringstyle=\color{codestr}, commentstyle=\color{codecmt}\itshape,
  columns=fullflexible, keepspaces=true, breaklines=true, breakatwhitespace=true,
  showstringspaces=false, tabsize=2, xleftmargin=4pt, xrightmargin=4pt,
  extendedchars=true, inputencoding=utf8,
}
\lstdefinelanguage{TypeScript}{
  keywords={const,let,function,return,if,else,for,while,of,in,new,export,import,from,type,interface,async,await,true,false,null,undefined,throw,class,extends,implements},
  sensitive=true, morecomment=[l]{//}, morecomment=[s]{/*}{*/}, morestring=[b]', morestring=[b]", morestring=[b]`,
}

% Bảng kẻ khung như mẫu: dòng cao hơn một chút
\renewcommand{\arraystretch}{1.25}
\setlength{\tabcolsep}{4.5pt}
\newcolumntype{L}[1]{>{\raggedright\arraybackslash}p{#1}}
\newcolumntype{C}[1]{>{\centering\arraybackslash}p{#1}}
\newcolumntype{R}[1]{>{\raggedleft\arraybackslash}p{#1}}

% Chỗ cần sinh viên bổ sung số liệu thật (không bịa)
\newcommand{\canbosung}[1]{\textcolor{red}{\textbf{[Cần bổ sung:} #1\textbf{]}}}

% Hình: \hinh{file}{độ rộng}{chú thích}{nhãn}
\newcommand{\hinh}[4]{%
  \begin{figure}[H]\centering\includegraphics[width=#2,height=0.8\textheight,keepaspectratio]{#1}\caption{#3}\label{#4}\end{figure}}

% Cho phép giãn chữ thêm để tránh dòng tràn lề với từ tiếng Việt dài
\setlength{\emergencystretch}{3em}
